\clearpage
\section{Exploratory Inventory-Management Extension}\label{sec:19}\label{app:i}
This appendix reports a separate exploratory joint-contract experiment. It is not part of the fitted Q-learning allocation test.

\subsection{Management Rule and Results}
Entry sizing and strategy allocation leave a third question: should an existing option position be trimmed as risk accumulates? A separate experiment uses a \$10 million joint whole-contract account, fixed 90\% put / 10\% straddle budgets and half-size entries in both sleeves. It does not replace the bandit accounts or the fractional-index Q experiment above.


A nonnegative regression predicts the magnitude of next-session straddle-book market P\&L from gross straddle vega and the preceding 21-session RMS P\&L. The coefficients and a 75th-percentile risk-per-NAV cutoff use 2018--2021 data. Cuts are inactive then, so this fits the warning without validating the management rule. When the reference-book warning is high, new straddles pause and existing cohorts target half their original integer counts at the next close. Trims remain in place; put rules are unchanged. This is an interpretable risk rule, not a learned Q-management policy.


\begin{figure}[!htbp]\centering
\includegraphics[width=\linewidth]{straddle_management.pdf}
\caption{Separate joint-contract management trial, January 2022--August 2025. Upper panel: compounded growth relative to cash. Lower panel: cumulative daily excess-return difference, conditional cuts minus unchanged 90/10. This is not the Q-index account.}\label{fig:straddle-management}
\end{figure}

\begin{table}[H]\centering\small
\caption{Separate joint-contract management trial, January 2022--August 2025.}
\setlength{\tabcolsep}{4pt}\renewcommand{\arraystretch}{1.18}
\begin{tabularx}{\textwidth}{@{}>{\hsize=1.87500\hsize\linewidth=\hsize\raggedright\arraybackslash}X>{\hsize=0.83008\hsize\linewidth=\hsize\raggedleft\arraybackslash}X>{\hsize=0.78125\hsize\linewidth=\hsize\raggedleft\arraybackslash}X>{\hsize=0.68359\hsize\linewidth=\hsize\raggedleft\arraybackslash}X>{\hsize=0.83008\hsize\linewidth=\hsize\raggedleft\arraybackslash}X@{}}\toprule
\textbf{Account} & \textbf{Mean excess (\%)} & \textbf{Vol. (\%)} & \textbf{Sharpe} & \textbf{Drawdown (\%)}\\\midrule
Unchanged 90/10 & 0.718 & 0.525 & 1.365 & 0.676\\
Conditional cuts & 0.733 & 0.465 & 1.576 & 0.734\\
Constant 90/5 + cash & 0.499 & 0.435 & 1.147 & 0.719\\
\bottomrule\end{tabularx}\end{table}

During January 2022--August 2025, conditional cuts raise excess Sharpe from 1.365 to 1.576, mainly by reducing volatility while preserving the arithmetic mean excess return. Cash-relative maximum drawdown becomes slightly worse, 0.734\% versus 0.676\%. The paired 126-session bootstrap interval for the Sharpe gain is -0.22 to +0.66. Annualized mean excess return is 252 times the daily mean, not compounded growth.


The constant-lower control assigns 90\% to puts, 5\% to straddles and 5\% to cash. It takes less average straddle exposure than the conditional policy, so it does not isolate timing at equal exposure. The trial motivates inventory-aware learning but does not establish that trimming reliably improves performance. Sections~\ref{app:trimming-protocol} and~\ref{app:trimming-scope} provide the timing, exposure and reproduction details.


\FloatBarrier
\FloatBarrier
\subsection{Risk Model, Execution, and Uncertainty}\label{app:trimming-protocol}


The warning model uses nonnegative coefficients 1.405401 on gross straddle vega and 0.383233 on 21-session RMS market P\&L. Divide its dollar prediction by prior NAV and compare with 0.0001616954. The predictor targets absolute next-session market P\&L, not its sign. A causally maintained unchanged reference book supplies the state, so a trim cannot mechanically switch off its own warning.


A high warning pauses new straddles and requests 50\% of each original cohort's call and put counts, rounded down. One-contract cohorts close both legs. The target is retried until filled or expired and never restores trimmed positions; new entries resume after the warning clears. The first signal is 31 December 2021, execution 3 January 2022, and first affected market return 4 January. All 908 warning records strictly order signal, execution and first affected return.


Mean later gross straddle vega is \$562.36 per volatility point for unchanged 90/10, \$450.59 for conditional cuts and \$275.54 for constant 90/5 plus cash. Thus the constant-lower control does not match average exposure. The saved paired bootstrap uses 2,000 draws; the Sharpe-gain interval versus holding is [-0.243, +0.697] at 63 sessions and [-0.218, +0.665] at 126. These intervals condition on the retained histories without research-selection adjustment.



\subsection{Experimental Scope and Reproducibility}\label{app:trimming-scope}
The three management accounts have identical histories before 2022, each with excess Sharpe 0.882. No cuts are applied during that period; it estimates the warning model and cutoff rather than validating the trimming benefit. Their later daily return and benchmark identities reconcile, including whole-contract exposure and funding in the saved ledgers. Other forecast, sizing and allocation variants were investigated without a replacement policy being retained. This one bounded trial uses already researched historical dates and is a separate extension, not a fresh holdout or additional evidence that the main Q learner manages individual contracts.
The attribution and management snapshots share the reproduction boundary described in Section~\ref{app:reproduction}.
\FloatBarrier
